\scanpage{321}
et on a
\[
(38)\qquad (-\unsure{\delta v|\delta\nu|\delta_0})B=ac+bc-\unsure{v|\nu}bd
\]
ce qui redonne $\uline{u}$ en termes de $a,b,c,d$, \uline{grâce à} la formule $\uline{\overline{\mathrm{IV}}}(48)$ (p. $453'$)
\[
(39)\qquad \uline{u}=\begin{pmatrix}\dfrac1D&\inserted{B/D}\deleted{\unsure{\underline u=}}\\0&1\end{pmatrix}
=\begin{pmatrix}
\dfrac{-\unsure{v|\nu}(ad-bc)}{\unsure{-vd^2+cd+d^2|-vc^2+cd+d^2}}&\dfrac{-ac+bc-\unsure{v|\nu}bd}{\unsure{-va^2+cd+d^2|-vc^2+cd+d^2}}\\0&1
\end{pmatrix}.
\]

\uline{Remarque}. Par contre, la donnée de $\beta$ \inserted{\uline{seule}} \inserted{(sans le ``signe'' $\varepsilon$)} ne permet pas de récupérer $\mu$. Mais la donnée de $(\beta,\unsure{\varepsilon})$ le permet, car on en déduit la connaissance de
\[
(\unsure{40})\qquad \eta=\varepsilon\beta\sigma(\beta)^{-1}\overset{\mathrm{dfn}}=\begin{pmatrix}R&S\\T&U\end{pmatrix},
\]
qui sera relié à $\zeta=\alpha\rho(\alpha)^{-1}$ par
\[
\zeta=\eta\varepsilon_1(\mu-1)
\]
et par la théorie de (II) on aura
\[
(41)\qquad \mu=1/U
\]
(mais en termes de $\beta\sigma(\beta)^{-1}\overset{\mathrm{dfn}}=\begin{pmatrix}R_0&S_0\\T_0&U_0\end{pmatrix}$,
on aura \deleted{\illegible{1}})
\[
(42)\qquad \mu=\frac{\varepsilon}{U_0},
\]
i.e. il n'est \inserted{a priori} pas connu par $U_0$ seul, mais seulement via $U_0$ et $\varepsilon$. \deleted{\illegible{1}}

Pour préciser ce point, on va regarder \inserted{quand et} comment, pour $(\beta,\varepsilon)$ fixés, on peut compléter par $\alpha,\mu$ de façon à avoir une
\mspage{320}

\scanpage{322}
\mspage{461}
section $(\alpha,\beta,\mu,\varepsilon)$ de $\unsure{\mathcal G^{*}_{\rho,\sigma}}$. Une condition nécessaire est que l'on ait $U_0$ inversible (elle ne dépend que de $\beta$, pas de $\varepsilon$) --- et alors on peut choisir $\varepsilon$ arbitrairement dans $\Gamma\mu_2$, \deleted{et}\inserted{$\varepsilon\in\Gamma\mu_2$} $\eta=\varepsilon\beta\sigma(\beta)^{-1}$ sera alors une section du \inserted{\unsure{ss-schéma}} $A_\sigma$ \deleted{\illegible{1}} de $G$ des solutions de $\det\eta=1$, $\operatorname{Tr}\eta\sigma=\operatorname{Tr}\sigma$ (cf lemme p. 455), donc par la théorie du (II), pour \deleted{\unsure{compléter prenant}} $\mu=1/U=\varepsilon/U_0$, \deleted{\illegible{2}} et définissant
\[
\zeta=\eta\varepsilon_1(\mu-1),
\]
$\zeta$ pourra bel et bien s'écrire \inserted{loc.t} sous la forme $\alpha\rho(\alpha)^{-1}$, l'indétermination étant d'ailleurs dans $T_\rho$ via $\alpha\longmapsto\alpha'=\alpha\uline b$, $\uline b\in\Gamma T_\rho$.

Ceci montre bien que $\mu$ \unsure{en|ne} peut prendre, pour une $\beta$ fixée, ttes les valeurs $\varepsilon/U_0$, où $\beta\sigma(\beta)^{-1}=\begin{pmatrix}R_0&\unsure{S_0|C_0}\\T_0&U_0\end{pmatrix}$ \ldots
\mspage{321}

\scanpage{323}
\mspage{462}
($\uline{\overline{\mathrm{VI}}}$) Caractérisation du triplet $(\varepsilon_1,\rho,\sigma)$ modulaire.

$\unsure{1^\circ})$ Soient $G,SG,PG$ des formes associées de $\mathrm{Gl}(2),\mathrm{Sl}(2),\mathrm{GP}(1)$ sur une base $S$ --- correspondant donc à : un torseur sous $\mathrm{GP}(1)_S$, ou encore à un fibré en droites projectives \inserted{$P$} sur $S$, ou une Algèbre d'Azumaya $M$ de rang $\unsure{4|1}$, et considérons deux sections \inserted{\illegible{1}} $\widetilde\rho,\widetilde\sigma$ de $PG$, telles que
\[
(1)\qquad \unsure{a)}\quad\widetilde\varepsilon_1=\widetilde\rho\widetilde\sigma^{-1}
\]
soit une section unipotente \uline{régulière}, et que la section
\[
(2)\qquad \widetilde\varepsilon_0\overset{\mathrm{dfn}}=\widetilde\rho^{-1}(\widetilde\varepsilon_1)=\widetilde\sigma^{-1}(\widetilde\varepsilon_1)
\]
qui s'en déduit soit \deleted{\illegible{2}} \inserted{en chaque fibre} inéquivalente à $\unsure{\widetilde\varepsilon_0|\widetilde\varepsilon_1}$ (i.e. que les sections $x_0,x_1$ de $P$ qui correspondent à $\varepsilon_0,\varepsilon_1$ soient \deleted{\illegible{1}} distinctes en chaque fibre).

\margin{\uline{NB} Si $\varepsilon_0,\varepsilon_1$ sont les relèvements canoniques \illegible{1} de $\widetilde\varepsilon_0,\widetilde\varepsilon_1$ dans $\unsure{G|P}$, \unsure{réguliers}, on peut relever $\widetilde\rho,\widetilde\sigma$ \unsure{de façon} \uline{unique} en $\rho,\sigma$ de façon à satisfaire
$-\sigma^{-1}\rho=\varepsilon_0\quad(\Longleftrightarrow-\rho\sigma^{-1}=\varepsilon_1)$,
$(\rho+\sigma)^2=(\rho+\sigma)$.
\inserted{avec $\unsure{\widetilde\sigma(\varepsilon_0)=\widetilde\rho(\varepsilon_0)=\varepsilon_1|\widetilde\sigma(\widetilde\varepsilon_0)=\widetilde\rho(\widetilde\varepsilon_0)=\widetilde\varepsilon_1}$}.}

Alors il existe un épinglage \uline{unique}$^*$
\[
(3)\qquad\deleted{\unsure{G_0=}}\quad\mathrm{Gl}(2)_S\xrightarrow[\sim]{\varphi}G
\]
et des sections $\unsure{v|\nu},\unsure{w|\omega},\unsure{v|\nu}',\unsure{w|\omega}'$ de $\mathcal O_S$, satisfaisant
\[
(4)\qquad \unsure{v|\nu}+\unsure{v|\nu}'=0,\qquad \unsure{w|\omega}+\unsure{w|\omega}'=1,\qquad \unsure{v|\nu}\in\Gamma\mathcal O_S^*
\]
telles que $\widetilde\rho,\widetilde\sigma$ soient les images, par le composé $\mathrm{Gl}(2)_S\xrightarrow[\sim]{\varphi}G\longrightarrow PG$, des sections
\[
(5)\qquad \rho=\begin{pmatrix}0&1\\\unsure{v|\nu}&\unsure{w|\omega}\end{pmatrix},\qquad
\sigma=\begin{pmatrix}0&-1\\\unsure{v|\nu}'&\unsure{w|\omega}'\end{pmatrix}
\]
\margin{\uline{NB}$^*$ L'unicité \inserted{de $\unsure{v|\nu},\unsure{w|\omega}$ (i.e. de $\rho_0,\sigma_0$)} \deleted{\unsure{n'est pas vraie}} demande quelques restrictions sur $S$, il suffit que \unsure{tte section} $\unsure{\tau}$ de $\mathcal O_S$ telle que $\unsure{2\tau=0,\ \tau^2=0}$ soit nulle.}
\mspage{322}

\scanpage{324}
de $\mathrm{Gl}(2)_S$ --- et inversement, pour tt épinglage $\varphi$ et \unsure{tout système} \deleted{\illegible{4}}\inserted{$\unsure{v|\nu},\unsure{w|\omega},\unsure{v|\nu}',\unsure{w|\omega}'$} satisfaisant (4) --- i.e.
\[
(4\mathrm{bis})\qquad \unsure{v|\nu}'=-\unsure{v|\nu},\quad \unsure{w|\omega}'=1-\unsure{w|\omega},\quad\begin{gathered}\unsure{v|\nu}\in\Gamma(S,\mathcal O_S^*)\\\unsure{w|\omega}\in\Gamma(S,\mathcal O_S)\end{gathered},
\]
les images \inserted{$\widetilde\rho,\widetilde\sigma$} par le composé $\mathrm{Gl}(2)_S\longrightarrow PG\simeq G/\mathbf G_m$ \deleted{satisfont}\inserted{des} sections $\rho,\sigma$ de $\mathrm{Gl}(2)_S$ données par (5) satisfont les conditions plus haut ($\widetilde\rho\widetilde\sigma^{-1}\overset{\mathrm{dfn}}=\widetilde\varepsilon_1$ est unipotent \uline{régulier}, et partout non équivalent à $\widetilde\sigma^{-1}(\widetilde\varepsilon_1)=\widetilde\rho^{-1}(\widetilde\varepsilon_1)\overset{\mathrm{dfn}}=\widetilde\varepsilon_0$).

Posant alors
\[
(6)\qquad \varepsilon_{\unsure{1|0}}=\begin{pmatrix}1&0\\1&1\end{pmatrix},\qquad \varepsilon_0=\begin{pmatrix}1&1/\unsure{v|\nu}\\0&1\end{pmatrix}
\]
on aura
\[
(7)\qquad \rho\sigma^{-1}=-\varepsilon_1,\qquad\sigma^{-1}\rho=-\varepsilon_0
\]
\[
(7\mathrm{bis})\qquad \varepsilon_0=\rho^{-1}(\varepsilon_1)=\sigma^{-1}(\varepsilon_1),\qquad\varepsilon_1=\rho(\varepsilon_0)=\sigma(\varepsilon_0)
\]
\uline{d'où} les images par $\mathrm{Gl}(2)_S\longrightarrow PG$ de \deleted{$\rho$ et $\sigma$} $\varepsilon_1,\varepsilon_0$ sont $\widetilde\varepsilon_1,\widetilde\varepsilon_0$. On a les formules
\[
(8)\qquad \operatorname{Tr}\sigma=\unsure{w|\omega}',\quad\operatorname{Tr}\rho=\unsure{w|\omega},\quad\operatorname{Tr}(\sigma+\rho)=1
\]
\[
(9)\qquad \sigma+\rho=\begin{pmatrix}0&0\\0&1\end{pmatrix}\overset{\mathrm{dfn}}=N\deleted{\unsure{0}}
\]
\[
(10)\qquad \operatorname{Tr}\zeta N=D\quad\text{si}\quad\zeta=\begin{pmatrix}A&B\\C&\uline D\end{pmatrix},
\]
\[
(11)\qquad \det\sigma=\det\rho=-\unsure{v|\nu}=\unsure{v'|-v'}.
\]
\mspage{323}

\scanpage{325}
\mspage{463}
$2^\circ)$ Les conditions suivantes sont équivalentes
\diagram{\[
(12)\quad\left\{
\begin{tikzcd}[ampersand replacement=\&,column sep=large,row sep=large]
\boxed{\operatorname{Tr}\sigma=0}\arrow[r,Leftrightarrow]\arrow[d,Leftrightarrow]\&\operatorname{Tr}\rho=1\arrow[r,Leftrightarrow]\arrow[d,Leftrightarrow]\&\boxed{\widetilde\sigma^2=1}\arrow[d,Leftrightarrow]\\
\boxed{\unsure{w|\omega}'=0}\&\boxed{\unsure{w|\omega}=1}\&\sigma^2\in\Gamma\mathbf G_m\arrow[d,Leftrightarrow]\\
\&\&\boxed{\sigma^2=\unsure{v|\nu}\,\mathrm{id}.}
\end{tikzcd}\right.
\]}
Elles équivalent encore à ceci : désignant par $N'_\sigma$ le sous-schéma en groupes de $\mathrm{Gl}(2)$ défini par la condition
\[
(13)\qquad \tau(\widetilde\sigma)=\widetilde\sigma\quad\text{i.e.}\quad\tau(\sigma)=\varepsilon\sigma,\quad\varepsilon\in\Gamma\mathbf G_m
\]
(on aura nécessairement $\varepsilon^2=1$ i.e. $\varepsilon\in\Gamma\mu_2$)
le morphisme
\[
(14)\qquad\begin{aligned}N'_\sigma&\longrightarrow\mu_2\\\tau&\longmapsto\varepsilon\end{aligned}
\]
est un \uline{épimorphisme} --- de façon qu'on aura une suite exacte
\[
(15)\qquad1\longrightarrow T_\sigma\longrightarrow N'_\sigma\longrightarrow\mu_2\longrightarrow1,
\]
où \deleted{\illegible{1}}
\[
(16)\qquad T_\sigma=\operatorname{Centr}_{\mathrm{Gl}(2)_S}(\sigma)=\{t\sigma+u\mid\det(t\sigma+u)\ \text{inv}\}
\]
\[
(17)\qquad\text{On a}\quad N'_\sigma\subset N_\sigma\overset{\mathrm{dfn}}=\operatorname{Norm}_{\mathrm{Gl}(2)}(T_\sigma).
\]
Si $2$ est section régulière de $\mathcal O_S$, alors tte section de $N_\sigma$ \inserted{sur $S$} est une section de $N'_\sigma$ --- donc \unsure{en car. rés. est $\ne2$} i.e. $2$ inv, on aura
\[
(18)\qquad N'_\sigma=N_\sigma\qquad(\text{car. rés. }\ne2).
\]
La condition envisagée signifie que $\sigma,\rho$ prennent la forme
\mspage{324}

\scanpage{326}
\[
(19)\qquad\sigma=\begin{pmatrix}0&-1\\\underset{\substack{\Vert\\-\unsure{v|\nu}}}{\unsure{v|\nu}'}&0\end{pmatrix},\qquad\rho=\begin{pmatrix}0&1\\\unsure{v|\nu}&1\end{pmatrix}
\]
avec $\unsure{v|\nu}+\unsure{v|\nu}'=0$.

Ces résultats de $1^\circ),2^\circ)$ ne sont qu'une récapitulation d'une partie des résultats des sections $\uline{\overline{\mathrm{IV}}},\uline{\overline{\mathrm{V}}}$ précédentes. \deleted{\illegible{1}} On vient d'étudier des conditions qui \unsure{reviennent|mènent} à la détermination \inserted{des valeurs de} $\unsure{w|\omega}$ (ou encore \inserted{de} $\unsure{w|\omega}'$) de façon bien explicite. On va \unsure{se tourner vers} des conditions qui vont servir à fixer $\unsure{v|\nu}$ (ou encore $\unsure{v|\nu}'$) de façon explicite.

$3^\circ)$ Conditions équivalentes
\[
(20)\quad\left\{\begin{aligned}
\det\sigma=1&\Longleftrightarrow\det\rho=1\Longleftrightarrow \unsure{v|\nu}=-1\Longleftrightarrow \unsure{v|\nu}'=1\\
&\Longleftrightarrow\varepsilon_0=\begin{pmatrix}1&-1\\0&1\end{pmatrix}.
\end{aligned}\right.
\]
Ces équivalences résultent aussitôt de (6) et (11): les conditions signifient aussi que $\rho$ et $\sigma$ ont la forme
\[
(21)\qquad \deleted{\rho}\inserted{\sigma}=\begin{pmatrix}0&-1\\1&\unsure{w|\omega}'\end{pmatrix},\qquad\rho=\begin{pmatrix}0&1\\-1&\unsure{w|\omega}\end{pmatrix}
\]
\[
(\unsure{w|\omega}+\unsure{w|\omega}'=1)
\]
$4^\circ)$ Conjonction de $2^\circ)$ et $3^\circ)$.

Conditions équivalentes :
\mspage{325}

\scanpage{327}
\mspage{464}

\[
(22)\quad
\left\{
\begin{array}{ccccc}
\begin{gathered}
\text{Conjonction des}\\
\text{conditions de }2^\circ\\
\text{et des conditions}\\
\text{de }3^\circ
\end{gathered}
&\Longleftrightarrow&
\begin{gathered}\boxed{\sigma^2=-1}\\
\text{i.e. }\sigma^2+1=0\end{gathered}
&\Longleftrightarrow&
\boxed{\widetilde\sigma^2=\widetilde\rho^3=1}
\\
\Updownarrow&&&&\\
\unsure{v|\nu}=-1,\unsure{w|\omega}=1&\Longleftrightarrow&
\rho=\begin{pmatrix}0&1\\-1&1\end{pmatrix}
&\Longleftrightarrow&\operatorname{Tr}\rho=\det\rho=1
\\
\Updownarrow&&&&\\
\unsure{v|\nu}'=1,\unsure{w|\omega}'=0&\Longleftrightarrow&
\sigma=\begin{pmatrix}0&-1\\1&0\end{pmatrix}
&\Longleftrightarrow&(\operatorname{Tr}\sigma=0,\det\sigma=1)
\end{array}\right.
\]
\margin{\diagram{\begin{tikzcd}[
ampersand replacement=\&,
column sep=large,
row sep=small,
execute at end picture={\draw[-] (quadratic.north east) -- (quadratic.north west) -- (quadratic.south west) -- (quadratic.south east);}
]
\boxed{\widetilde\sigma^2=\widetilde\rho^3=1}
\arrow[r,Leftrightarrow]
\arrow[dr,Leftrightarrow]
\& \deleted{\boxed{\rho^3=-1}}
\\
\& |[alias=quadratic,inner sep=3pt]| {\rho^2-\rho+1=\unsure{\illegible{1}}}
\end{tikzcd}}
cf aussi plus bas}

\uline{Dém.} Le fait que la conjonction des conditions de $2^\circ)$ et $3^\circ)$ soit équivalente à \unsure{chacune} des conditions des $2^{\mathrm e}$ et $3^{\mathrm e}$ lignes, est trivial. L'équivalence avec $\sigma^2=-1$ provient du fait que \deleted{\illegible{1}} la condition $2^\circ$ signifie que \inserted{$\sigma^2$ est scalaire, soit} $\sigma^2=\unsure{v|\nu}\,\mathrm{id}$, \deleted{\illegible{2}} tandis que $3^\circ$ signifie $\unsure{v|\nu}=-1$, donc $\sigma^2=-1$ équivaut bien à la conjonction.

\margin{pour une autre démonstration de l'équivalence de $2^\circ)+3^\circ)$ avec $\widetilde\sigma^2=\widetilde\rho^3=1$, cf lemme ci-dessous}

Notons que $2^\circ)$ s'écrit aussi $\widetilde\sigma^2=1$, et équivaut à $\unsure{w|\omega}'=0$, soit $\unsure{w|\omega}=1$, et on trouve alors, en écrivant \deleted{\unsure{$\rho=$}} 
\[
\rho=\begin{pmatrix}0&1\\\unsure{v|\nu}&1\end{pmatrix},
\]
\[
(23)\qquad
\rho^3=\begin{pmatrix}\unsure{v|\nu}&\unsure{v|\nu}+1\\\unsure{v|\nu}(\unsure{v|\nu}+1)&2\unsure{v|\nu}+1\end{pmatrix}
\]
donc $\rho^3$ ne peut être un scalaire i.e. $\widetilde\rho^3=1$ que si les deux termes diagonaux sont égaux, or leur diff est $(2\unsure{v|\nu}+1)-\unsure{v|\nu}=\unsure{v|\nu}+1$,

\mspage{326}

\scanpage{328}

donc la condition s'écrit $\unsure{v|\nu}=-1$, i.e. c'est la condition $3^\circ$, mais alors l'expression (23) de $\rho^3$ nous montre que $\rho^3=-1$.

\uline{\unsure{Remarques|Remarque}} Les conditions qui précèdent impliquent les formules
\[
\begin{aligned}
(24)\qquad &\varepsilon_0=\sigma\rho\\
(25)\qquad &\varepsilon_1=\rho\sigma\\
(26)\qquad &\rho\varepsilon_1=\varepsilon_0^{-1}\\
(27)\qquad &\rho^2\sigma=\varepsilon_0^{-1}
\end{aligned}
\]
dont chacune \inserted{sauf la dernière (27)} les implique à son tour
--- p.ex. on a $\varepsilon_0=-\sigma^{-1}\rho$ \inserted{(7)}, donc (24) équivaut à \inserted{$\sigma=-\sigma^{-1}$ i.e.} $\sigma^2=-1$, et (25) équivaut à (24) en appliquant $\operatorname{int}(\rho)$ aux deux \unsure{mbs}.

On a d'ailleurs
\[
\rho\varepsilon_1=
\begin{pmatrix}0&1\\\unsure{v|\nu}&\unsure{w|\omega}\end{pmatrix}
\begin{pmatrix}1&0\\1&1\end{pmatrix}
=\begin{pmatrix}1&1\\\unsure{v|\nu}+\unsure{w|\omega}&\unsure{\unsure{w|\omega}^{\prime}|\unsure{w|\omega}}\end{pmatrix},
\qquad
\varepsilon_0^{-1}=\begin{pmatrix}1&1\\0&1\end{pmatrix}
\]
qui montre que $\rho\varepsilon_1=\varepsilon_0^{-1}$ équivaut à :
\[
\unsure{w|\omega}=1,\quad \unsure{v|\nu}+\unsure{w|\omega}=0\qquad\text{i.e.}\quad \unsure{w|\omega}=1,\unsure{v|\nu}=-1\quad\text{i.e.}\quad
\rho=\begin{pmatrix}0&1\\-1&1\end{pmatrix},
\]
OK. Composant (25) \inserted{(sous la forme $\rho\varepsilon_1=\rho^2\sigma$)} et (26) on trouve (27) comme conséquence des conditions $2^\circ)+3^\circ)$ ; \deleted{\unsure{l'}}\inserted{Essayons de prouver} implication inverse \deleted{\unsure{résulte}}\inserted{par} \deleted{d'}un calcul brutal
\[\deleted{\rho^2\sigma=}o^2\sigma=
\begin{pmatrix}
-\unsure{v|\nu}\unsure{w|\omega}&\unsure{w|\omega}-(\unsure{v|\nu}+\unsure{w|\omega}^2)\\
-\unsure{v|\nu}(\unsure{v|\nu}+\unsure{w|\omega}^2)&\unsure{-\unsure{w|\omega}(\unsure{w|\omega}^2+2\unsure{v|\nu})+\unsure{v|\nu}-\unsure{w|\omega}^2|-\unsure{w|\omega}(\unsure{w|\omega}^2+2\unsure{v|\nu})+\unsure{v|\nu}+\unsure{w|\omega}^2}
\end{pmatrix},
\qquad
\varepsilon_0^{-1}=\begin{pmatrix}1&\unsure{-1/\unsure{v|\nu}}\\0&1\end{pmatrix}
\]
\deleted{comme $\unsure{v|\nu}$ inversible} on ne peut avoir l'égalité $\rho^2\sigma=\unsure{-\varepsilon_0^{-1}|\varepsilon_0^{-1}}$ que si $-\unsure{v|\nu}(\unsure{v|\nu}+\unsure{w|\omega}^2)=0$, \unsure{cela implique|ce qui implique} ($\unsure{v|\nu}$ inv. !) $\unsure{v|\nu}+\unsure{w|\omega}^2=0$, \deleted{\unsure{d'où}}\inserted{i.e. $\unsure{v|\nu}=-\unsure{w|\omega}^2$},

\mspage{327}

\scanpage{329}
\mspage{465}

\[
\rho^2\sigma=\begin{pmatrix}\unsure{w|\omega}^3&\unsure{w|\omega}\\0&\deleted{\unsure{-w(1+2w)}}\inserted{\unsure{w|\omega}^3}\end{pmatrix},
\]
Cela montre que
\deleted{$\rho^2\sigma=\varepsilon_0^{-1}$ \unsure{ssi} $\unsure{w|\omega}=1$ OK.}
\margin{\deleted{\unsure{On a en plus de} \illegible{4}}
\unsure{$\unsure{v|\nu}=-\unsure{w|\omega}^2$} \deleted{\illegible{2}}
\unsure{i.e. $\unsure{v|\nu}=-1,\unsure{w|\omega}=1$}}
\inserted{
\[
(27\mathrm{bis})\qquad
\rho^2\sigma=\varepsilon_0^{-1}\Longleftrightarrow
\left\{\unsure{v|\nu}+\unsure{w|\omega}^2=0,\ \unsure{w|\omega}^3=1\right.
\quad
\left(\Rightarrow \unsure{v|\nu}=-\unsure{w|\omega}^2=-1/\unsure{w|\omega}\ \text{i.e.}\ \unsure{w|\omega}=-1/\unsure{v|\nu},\ \mathrm{OK}\right)
\]
}

Mais (27), en multipliant par $\rho$, équivaut aussi à
\[\deleted{\unsure{(27)}\,\illegible{1}}ho^3\sigma=\deleted{\unsure{u}}\rho\varepsilon_0^{-1}
\]
or on sait que $\varepsilon_0=\unsure{-\sigma^{-1}\rho}$ \inserted{(7)} i.e. $\varepsilon_0^{-1}=-\deleted{\unsure{\sigma}}\rho^{-1}\unsure{\sigma|\sigma_0}$
i.e. $\rho\varepsilon_0^{-1}=-\sigma$, donc la relation précédente s'écrit aussi $\rho^3\sigma=-\sigma$ ou encore
\[
(28)\qquad \boxed{\rho^3=-1}.
\]
On peut voir \deleted{\unsure{cette}} l'équivalence \deleted{directement}\inserted{de (28) avec les conditions (27bis)} en calculant $\rho^3$ en fonction de $\unsure{v|\nu},\unsure{w|\omega}$, on trouve
\[
(29)\qquad
\rho^3=\begin{pmatrix}\unsure{v|\nu}\unsure{w|\omega}&\unsure{v|\nu}+\unsure{w|\omega}^2\\\unsure{v|\nu}(\unsure{v|\nu}+\unsure{w|\omega}^2)&\unsure{w|\omega}(2\unsure{v|\nu}+\unsure{w|\omega}^2)\end{pmatrix}
\]
Et on ne peut avoir $\widetilde\rho^3=1$ i.e. $\rho^3$ un scalaire que si \deleted{$\unsure{v|\nu}+\unsure{w|\omega}^2=\illegible{1}$}
\[
\unsure{v|\nu}+\unsure{w|\omega}^2=0\qquad\text{i.e.}\quad \unsure{v|\nu}=-\unsure{w|\omega}^2
\]
ce qui implique
\[
\rho^3=\begin{pmatrix}-\unsure{w|\omega}^3&0\\0&-\unsure{w|\omega}^3\end{pmatrix}
\]
i.e. on trouve

\uline{Lemme} On a $\widetilde\rho^3=1$ ssi $\unsure{v|\nu}+\unsure{w|\omega}^2=0$ i.e. $\unsure{v|\nu}=-\unsure{w|\omega}^2$ i.e. on a
\[
(30)\qquad
\rho=\begin{pmatrix}0&1\\-\unsure{w|\omega}^2&\unsure{w|\omega}\end{pmatrix},
\]
et alors on a
\[
(31)\qquad \rho^3=\unsure{-}\unsure{w|\omega}^3\cdot1\ ,
\]
donc on a
\[
(32)\qquad
\rho^3=-1\Longleftrightarrow
\deleted{\unsure{\unsure{v|\nu}=-\unsure{w|\omega}^2}}\quad
\begin{gathered}
\unsure{w|\omega}^3=1\ \text{i.e. }\unsure{w|\omega}\in\Gamma\mu_3\\
\inserted{\uline{et}\quad \unsure{v|\nu}=-\unsure{w|\omega}^2\ (= -1/\unsure{w|\omega})}
\end{gathered}
\]

\mspage{328}

\scanpage{330}

\uline{Corollaire 1} Les conditions équivalentes (22) équivalent aussi à
\[
(33)\qquad \rho^3=-1\quad\text{et}\quad\det\rho=1
\]
ou encore à
\[
(34)\qquad \rho^3=-1\quad\text{et}\quad\deleted{\unsure{\det}}\inserted{\operatorname{Tr}}\rho=1
\]
en d'autres termes, considérant les trois relations
\[
(35)\qquad \det\rho=1,\quad\operatorname{Tr}\rho=1,\quad\rho^3=-1
\]
\inserted{la conjonction de} deux de celles-ci \unsure{impliquent|implique} la troisième, et équivaut aux conditions équivalentes de (22).

\margin{faux}
\deleted{
Notons une autre façon de trouver ce corollaire, en notant que
\[
(\unsure{36|30})\qquad
\rho^3+1=(\rho+1)(\rho^2-\rho+1)\qquad
\text{(\uline{identité} algébrique)}
\]
or $\rho$ n'étant un scalaire sur aucune fibre,
\uline{$\rho+1$ est inversible}, donc $\rho^3=-1$ i.e. $\rho^3+1=0$ équivaut à la relation $\rho^2-\rho+1=0$ :
\[
(36)\qquad \rho^2-\rho+1=0\ ;\quad(\Longleftrightarrow\rho^3=-1),
\]
qu'on compare avec Hamilton--Cayley
\[
(37)\qquad
\rho^2-\underbrace{(\operatorname{Tr}\rho)}_{\unsure{w|\omega}}\rho+\underbrace{\det\rho}_{-\unsure{v|\nu}}=0,
\]
donc \deleted{$\unsure{w|\omega}=1,\unsure{v|\nu}=-1$} $\operatorname{Tr}\rho=\det\rho=1$ implique (36), \deleted{donc $\rho^3=-1$}
d'autre part si on a (36), on conclut via 37
\[
(\operatorname{Tr}(\rho)-1)\rho\ \unsure{=|\ne}\ \deleted{\unsure{\det}}\,(\det\rho-1)1
\]
}

\mspage{329}

\scanpage{331}
\mspage{466}

Une autre façon de trouver certaines équivalences est via les équations de Hamilton--Cayley pour $\sigma$ et $\rho$
\[
(36)\qquad
\deleted{\unsure{\rho}}\inserted{\sigma}^2
-\underbrace{(\operatorname{Tr}\sigma)}_{\unsure{w|\omega}'=1-\unsure{w|\omega}}\sigma
+\underbrace{\det\sigma}_{\unsure{v|\nu}'=-\unsure{v|\nu}}=0
\]
\[
(37)\qquad
\rho^2-\underbrace{(\operatorname{Tr}\rho)}_\unsure{w|\omega}\rho
+\underbrace{\det\rho}_{-\unsure{v|\nu}}=0.
\]
L'équation (36) redonne que la conjonction de $\operatorname{Tr}\sigma=0$, $\det\sigma=1$ implique $\sigma^2+1=0$ i.e. $\sigma^2=-1$, et inversement, compte tenu que $1$ et $\sigma$ sont \inserted{des sections} \unsure{lin\textsuperscript{t}} indépendantes \unsure{en|de} $\unsure{\mathrm{M}_2}(\unsure{\underline{O}_S|\mathcal O_S})$.
De $\unsure{\overline{\mathrm m}|=}$ (37) montre
\[
(38)\qquad
(\operatorname{Tr}\rho=1\ \text{et}\ \det\rho=1)
\Longleftrightarrow(\rho^2-\rho+1=0).
\]
et l'on \deleted{relation} a l'implication
\[
(39)\qquad
\rho^2-\rho+1=0\quad\Longrightarrow\quad
\rho^3+1=0\ \text{i.e.}\ \rho^3=-1
\]
en vertu de l'identité \uline{algébrique}
\[
\rho^3+1=(\rho+1)(\rho^2-\rho+1).
\]
Donc

\uline{Corollaire 2}. Les conditions équivalentes de (22) équivalent aussi à la relation
\[
(40)\qquad \rho^2+\rho+1=0.
\]

\mspage{330}

\scanpage{332}

\deleted{
$5^\circ)$ Théorème d'épinglage \unsure{fondateur|fondamental}.

\uline{Théorème} Soient $G,SG,PG$ des formes \inserted{(\unsure{sur schéma} $S$)} de $\mathrm{Gl}(2),\mathrm{Sl}(2),\mathrm{GP}(1)$ \inserted{\deleted{\unsure{une algèbre}}} \unsure{(}associées\unsure{)}, et $\widetilde\sigma,\widetilde\rho$ deux sections de $PG$.
}

$5^\circ)$ \uline{Un retour sur} $\uline{1^\circ)}$, relatif à la situation
\[
(41)\quad
\left\{\begin{array}{l}
\widetilde\rho,\widetilde\sigma\ \text{avec }1^\circ)\quad
\widetilde\varepsilon_0=\widetilde\sigma^{-1}\widetilde\rho\quad
\text{unipotent }\uline{\text{régulier}}\ \deleted{\text{\unsure{i.e.}}}\quad(\text{i.e.}\\
\widetilde\varepsilon_1=\widetilde\rho\widetilde\sigma^{-1}
=\widetilde\rho(\widetilde\varepsilon_0)
=\widetilde\sigma(\widetilde\varepsilon_0)\quad
\text{unipotent }\uline{\text{régulier}})\ \text{et}\\
2^\circ)\quad\widetilde\varepsilon_0\ \text{et}\ \widetilde\varepsilon_1\ \text{non équivalents en chaque fibre.}
\end{array}\right.
\]

Soient $\varepsilon_0,\varepsilon_1$ les relèvements unipotents \uline{réguliers} dans $G$ --- identifiés à des sections de $\unsure{M|\mathcal M}$
\[
(42)\qquad\varepsilon_0=1+N_0,\quad\varepsilon_1=1+N_1,
\]
\margin{$N_0$ et $N_1$ strictement \unsure{nilp.\textsuperscript{ts}} i.e. leur det et leur tr sont nuls}
avec
\[
(43)\qquad
\widetilde\sigma(\varepsilon_0)=\widetilde\rho(\varepsilon_0)=\varepsilon_1,
\quad\inserted{\text{i.e.,}}\quad
\widetilde\sigma(N_0)=\widetilde\rho(N_0)=N_1.
\]
On peut alors relever $\widetilde\sigma$ et $\widetilde\rho$ de façon canonique en des sections $\sigma,\rho$ de $G$ (ou de $\unsure{M|\mathcal M}$) satisfaisant les conditions suivantes
\[
(44)\quad\left\{
\begin{array}{ll}
1^\circ)\ -\sigma^{-1}\rho=\varepsilon_0
&[\Longleftrightarrow-\rho\sigma^{-1}=\varepsilon_1]\\
\deleted{\unsure{2^\circ)(\rho+\sigma)}}
&\inserted{\left[\begin{aligned}
&\Longleftrightarrow\rho+\sigma=-N_1\sigma\ (= -\sigma N_0)\\
&\Longleftrightarrow\rho+\sigma=N_1\rho\ (=\rho N_0)
\end{aligned}\right]}\\
&\deleted{\unsure{\Longleftrightarrow\rho+\sigma=}}\\
2^\circ)\ (\rho+\sigma)^2=\rho+\sigma
&\inserted{\begin{gathered}
\Longrightarrow(\rho+\sigma)^2=\unsure{u}(\rho+\sigma)\quad\text{avec}\\
u\in\Gamma G_m
\end{gathered}}
\end{array}\right.
\]
De plus, il existe un \deleted{\unsure{iso}} unique triplet $(\varphi,\rho_0,\sigma_0)$
\[
(45)\quad\left\{\begin{array}{l}
G_0=\mathrm{Gl}(2)_S\xrightarrow{\varphi}G\quad(\text{épinglage})\\
\rho_0,\sigma_0\in\Gamma(G_0),\quad
\rho_0=\begin{pmatrix}0&1\\\unsure{v|\nu}&\unsure{w|\omega}\end{pmatrix},\quad
\sigma_0=\begin{pmatrix}0&-1\\\unsure{v|\nu}'&\unsure{w|\omega}'\end{pmatrix}\\
(\unsure{v|\nu}+\unsure{v|\nu}'=0,\unsure{w|\omega}+\unsure{w|\omega}'=1),\quad \unsure{v|\nu}\in\Gamma G_m\\
\inserted{\text{i.e. }\unsure{v|\nu}'=-\unsure{v|\nu},\unsure{w|\omega}'=1-\unsure{w|\omega}}
\end{array}\right.
\]
tel qu'on ait
\[
(\unsure{46|48})\qquad\varphi(\rho_0)=\rho,\quad\varphi(\sigma_0)=\sigma,
\]
ce qui implique
\[
(47)\qquad\varphi((\varepsilon_0)_0)=\varepsilon_0,\quad
\varphi((\varepsilon_1)_0)=\varepsilon_1
\]

\mspage{331}

\scanpage{333}
\mspage{467}

où \deleted{\illegible{1}}
\[
(48)\qquad
\unsure{(\varepsilon_0)_0}=-\sigma_0^{-1}\rho_0=
\begin{pmatrix}1&\dfrac{1}{\unsure{v|\nu}}\\0&1\end{pmatrix},\qquad
\deleted{\unsure{(\varepsilon_1)_0}}
\inserted{(\varepsilon_1)_0=-\rho_0\sigma_0^{-1}}=
\begin{pmatrix}1&0\\1&1\end{pmatrix}
\]
--- et inversement, pour $\rho_0,\sigma_0$ comme dessus, et tout iso.
$\varphi:G_0\simeq\mathrm{Gl}(2)_S\xrightarrow{\sim}G$, les images
$\widetilde\rho,\widetilde\sigma$ de $\rho,\sigma$ dans $PG$ satisfont
$(41)$, et les images $\rho,\sigma$ de $\rho_0,\sigma_0$ par $\varphi$
sont des relèvements de $\widetilde\rho,\widetilde\sigma$ satisfaisant $(44)$.

\uline{Théorème d'épinglage modulaire.} Soient $G,SG,PG$ des formes
(\uuline{sur} schéma $S$) de $\mathrm{Gl}(2),\mathrm{Sl}(2),\mathrm{GP}(1)$,
\unsure{associées} à une \unsure{alg.} d'Azumaya $M$ de rg. $4$.

$1.)$ Appelons \uline{\deleted{\illegible{1}}\inserted{triplet} modulaire}
dans $PG$ un triplet de sections $(\widetilde\rho,\widetilde\sigma,
\widetilde\varepsilon_0)$ de $PG$, satisfaisant les conditions $(49)(50)(51)$.

\deleted{$(\unsure{49})\quad \widetilde\rho^3=\widetilde\sigma^2=1$, $\widetilde\varepsilon_0$ est unipotent \unsure{régulier}.}

\[
(49)\qquad \widetilde\varepsilon_0=\widetilde\sigma^{-1}\widetilde\rho
\quad\text{i.e.}\quad
\widetilde\rho^{-1}\widetilde\sigma\widetilde\varepsilon_0=1
\]
[de sorte que
\[
(49\ \mathrm{bis})\qquad
\widetilde\varepsilon_1\overset{\mathrm{dfn}}=
\widetilde\rho\widetilde\sigma^{-1}=
\widetilde\rho(\widetilde\varepsilon_0)=
\widetilde\sigma(\widetilde\varepsilon_0)
\]
];

$(50)\qquad \widetilde\rho^3=\widetilde\sigma^2=1$,\quad
$\widetilde\varepsilon_0$ est unipotent régulier
($\Longleftrightarrow\widetilde\varepsilon_1$ est unipotent \uline{régulier});

$(51)\qquad \widetilde\varepsilon_0$ et $\widetilde\varepsilon_1$
\deleted{\unsure{partout}} non équivalents \deleted{\unsure{sur les}}
\inserted{en chaque} fibre
(ce qui signifie, en introduisant les relèvements
\[
\varepsilon_0=1+N_0,\qquad \varepsilon_1=1+N_1
\]
\margin{($\unsure{N_0\text{ et }N_1}$ strictement nilpotents,\ldots\ cf $(42)$)}
à $G$, que $N_0$ et $N_1$ non homothétiques en tte fibre).

$2)$ Appelons triplet modulaire \unsure{de|dans} $G$ un
\mspage{332}

\scanpage{334}

triplet de sections $(\rho,\sigma,\varepsilon_0)$ satisfaisant les
\deleted{\unsure{cinq}} conditions suivantes

\[
\begin{aligned}
(52)\qquad &-\sigma^{-1}\rho=\varepsilon_0
\quad(\Longleftrightarrow\ \rho^{-1}\sigma\varepsilon_0=-1.)\\
&\left[\begin{aligned}
&\Longleftrightarrow\ -\rho\sigma^{-1}=\varepsilon_1
&&\text{où }\varepsilon_1=\rho\varepsilon_0\rho^{-1}\\
&\Longleftrightarrow\ -\rho\sigma^{-1}=\varepsilon_1
&&\text{où }\varepsilon_1=\sigma\varepsilon_0\sigma^{-1}
\end{aligned}\right]
\end{aligned}
\]
\[
(53)\qquad(\rho+\sigma)^2=\rho+\sigma
\]

$(54)\qquad\varepsilon_0$
\deleted{et $\varepsilon_1\overset{\mathrm{dfn}}=\rho\varepsilon_0\rho^{-1}=\sigma\varepsilon_0\sigma^{-1}$ \unsure{sont}}
\inserted{est unipotent régulier i.e. $\varepsilon_0=1+N_0$ avec}
$N_0$ \deleted{\unsure{unip.}} nilpotent régulier
($\mathrm{Tr}\,N_0=\det N_0=0$, $N_0$ non nul en chaque fibre).

\margin{$\Longleftrightarrow\left\{\begin{array}{l}
\det\varepsilon_0=1\\
\mathrm{Tr}\,\varepsilon_0=2\\
\varepsilon_0\ne1\text{ en tte fibre}
\end{array}\right.$}

($\Longleftrightarrow\ \deleted{\unsure{\varepsilon_0}}\ \varepsilon_1
\overset{\mathrm{dfn}}=\rho\varepsilon_0\rho^{-1}
=\sigma\varepsilon_0\sigma^{-1}$ est unipot. rég.)

$(55)\qquad\deleted{\rho}\cdot\varepsilon_0$ et $\varepsilon_1$ sont
inéquivalents en chaque fibre
($\Longleftrightarrow\ N_0$ et $N_1$ non homoth. en tte fibre).

$(56)\qquad$ On a une des quatre conditions équivalentes
\[
\begin{array}{cc}
\boxed{\sigma^2+1=0}&\boxed{\rho^2\unsure{+|-}\rho+1=0}\\[6pt]
\boxed{\mathrm{Tr}\,\sigma=0,\det\sigma=1}&
\boxed{\mathrm{Tr}\,\rho=\det\rho=1}
\end{array}
\quad.
\]

\deleted{\unsure{3}} $3)$ Appelons \uline{triplet modulaire}
\uline{\unsure{standard}} dans $\mathrm{Gl}(2)_S$
(resp. dans $\mathrm{GP}(1)_S$) le triplet
\[
(57)\quad\left(\rho_0=\begin{pmatrix}0&1\\-1&1\end{pmatrix},\quad
\sigma_0=\begin{pmatrix}0&-1\\1&0\end{pmatrix},\quad
\varepsilon_0=\begin{pmatrix}1&-1\\0&1\end{pmatrix}\right)
\]
(resp. son image $(\widetilde\rho_0,\widetilde\sigma_0,
\unsure{\widetilde\varepsilon_0})$ dans $\Gamma\mathrm{GP}(1)_S$) ---
qui est bel et bien un triplet modulaire.

Ceci posé, on a \deleted{des}\inserted{le diagramme \unsure{comm.} d'}applications
canoniques
\mspage{333}

\scanpage{335}
\mspage{468}

\[
\varphi\longmapsto(\varphi(\rho_0),\varphi(\sigma_0),\varphi(\varepsilon_0))
\]

$(58)$
\diagram{\begin{tikzcd}[ampersand replacement=\&,column sep=large,row sep=large]
\operatorname{Isom}(\mathrm{Gl}(2),G)
\arrow[r]\arrow[d,phantom,"\unsure{\wr}"]\&
\mathrm{TripMod}(G)\arrow[d]\&
(\rho,\sigma,\varepsilon_0)\arrow[d,mapsto]\\
\operatorname{Isom}(\mathrm{GP}(1),PG)\arrow[r]\&
\mathrm{Tripmod}(PG)\&
(\widetilde\rho,\widetilde\sigma,\widetilde\varepsilon_0)
\end{tikzcd}}
\[
\psi\longmapsto(\psi(\widetilde\rho_0),\psi(\widetilde\sigma_0),
\psi(\widetilde\varepsilon_0))
\]

et on a le

\uline{Théorème d'épinglage.} Toutes les flèches de $(58)$ sont bijections.

Ce théorème ne fait que résumer \deleted{les}\inserted{certains}
\deleted{\unsure{réflexions}}\inserted{des} résultats principaux des réflexions
\unsure{du|des} présent \S\ \uline{$\overline{\mathrm{VI}}$}.
Notons que l'injectivité des flèches horizontales signifie \inserted{aussi}
qu'un \deleted{\unsure{isom.}}
$\mathrm{Gl}(2)_S=G_0\longrightarrow G$
(resp. $\mathrm{GP}(1)_S=PG_0\longrightarrow\unsure{G|PG}$
est connu quand on connaît son effet sur $\sigma$ et $\rho$
(donc aussi sur $\varepsilon_0$)
\unsure{$\leftarrow$} resp. quand on connaît son effet sur
$\widetilde\sigma$ et $\widetilde\rho$) --- ce qui signifie aussi
la $\unsure{\overline{\mathrm m}|=}$ chose pour les automorphismes de
$\mathrm{Gl}(2)_S$ (resp. de $\mathrm{GP}(1)_S$)
ce qui signifie aussi que
\[
(59)\qquad T_\rho\cap T_\sigma=\deleted{\{1\}}\ G_m
\]
(l'intersection des centralisateurs de $\rho,\sigma$ dans
$\mathrm{Gl}(2)_S$ est \unsure{réduite|réduit} au centre $G_m$) ---
\unsure{resp.} la relation plus forte
\[
(60)\qquad N'_\rho\cap N'_\sigma=G_m,
\]
\mspage{334}

\scanpage{336}

où $N'_\rho$, $N'_\sigma$ sont les \unsure{ss-schémas} en groupes de $\mathrm{Gl}(2)_S$ définis \deleted{par} \inserted{\unsure{resp\textsuperscript{t}}} comme stabilisateurs de $\widetilde\rho$, $\widetilde\sigma$ dans $\unsure{\mathrm{Gl}(2)_S}$, i.e.
\[
(61)\qquad
\left\{\begin{aligned}
N'_\rho&=\{g\mid g(\rho)=\varepsilon\rho,\ \varepsilon\text{ scalaire}\},\\
N'_\sigma&=\{g\mid g(\sigma)=\varepsilon\sigma,\ \varepsilon\text{ scalaire}\}
\end{aligned}\right.
\]
\margin{NB on \unsure{aura} \unsure{néc\textsuperscript{t}} $\varepsilon^2=1$}

\uline{Rmq} Il est plausible, \deleted{que \illegible{5}} \inserted{\unsure{Vu ces résultats,}} que $\mathrm{GP}(1)_{\mathbb Z}$ --- en tant que faisceau en groupes ``absolu'' (sur $\operatorname{Spec}\mathbb Z$) pour fppf, est engendré par les hom.
\[
\begin{aligned}
\mathbb Z/2\mathbb Z&\xrightarrow{\alpha\,\mapsto\,\widetilde\sigma_0^{\alpha}}\mathrm{GP}(1)_{\mathbb Z}\\
\mathbb Z/3\mathbb Z&\xrightarrow{\beta\,\mapsto\,\widetilde\rho_0^{\beta}}\mathrm{GP}(1)_{\mathbb Z}\\
\mathbb G_a&\xrightarrow{\lambda\,\mapsto\,\widetilde\varepsilon_0(\lambda)}\mathrm{GP}(1)_{\mathbb Z}
\end{aligned}
\]
satisfaisant \unsure{la|les} \inserted{seule} relation (\unsure{R})
\[
(62)\qquad \widetilde\rho_0^{-1}\widetilde\sigma_0\widetilde\varepsilon_0(1)=1
\qquad\deleted{\unsure{(?)}}
\]
--- cela vaudrait bien la peine de tirer cette question au clair, ainsi que la question analogue pour $\mathrm{Gl}(2)_{\mathbb Z}$, où les \deleted{relations} groupes $\mathbb Z/2\mathbb Z$, $\mathbb Z/3\mathbb Z$ doivent \unsure{se remplacer} par $\mathbb Z/4\mathbb Z$, $\mathbb Z/6\mathbb Z$, et la relation (62), par
\[
(63)\qquad \rho_0^{-1}\sigma_0\widetilde\varepsilon_0(1)=\rho_0^3=\sigma_0^2.
\]
\mspage{335}

\scanpage{337}
\mspage{469}

\uline{Rmq} J'ai la nette impression que le paquet de conditions (51) à (56) est surabondant, et me demande p.ex. si \unsure{l'ens.|l'une} des trois conditions algébriques simples
\[
(64)\qquad \sigma^2+1=\rho^2-\rho+1=0,\qquad (\rho+\sigma)^2=\rho+\sigma
\]
\unsure{n'implique} pas les autres conditions \inserted{(en définissant $\varepsilon_0$ par (\unsure{51|52}))}. La troisième condition (64) signifie que
\[
(65)\qquad N=\rho+\sigma
\]
est \unsure{idempotent}
\[
(66)\qquad N^2=N,
\]
et en termes de $N$ on aura
\[
\sigma^{-1}N=\sigma^{-1}\rho+1\quad\text{i.e.}
\]
\[
(67)\qquad
\left\{\begin{aligned}
\varepsilon_0&\overset{\mathrm{dfn}}{=}-\sigma^{-1}\rho=1-\sigma^{-1}N=1+N_0
&\qquad N_0&=-\sigma^{-1}N\\[-2pt]
&&&=\sigma N\\
\varepsilon_1&\overset{\mathrm{dfn}}{=}-\rho\sigma^{-1}=1-N\sigma^{-1}=1+N_1
& N_1&=-N\sigma^{-1}=N\sigma
\end{aligned}\right.
\]
\margin{\uline{NB} Il faut \deleted{\unsure{peut-être}} \uline{préciser} \deleted{\unsure{que dans}} \illegible{4} \unsure{premières} relations \unsure{(64)} \illegible{3}
$\operatorname{Tr}\sigma=0$, $\det\sigma=1$, $\operatorname{Tr}\rho=\unsure{0|1}$, $\det\rho=1$ \unsure{doit exiger que} \illegible{2} sont \uline{réguliers} i.e. \unsure{non homothéties} en aucune fibre \illegible{1}.}

et la condition $\varepsilon_0$ \deleted{\unsure{strict\textsuperscript{t} unipotent}} \inserted{\unsure{unipotent régulier}} s'exprime par les conditions
\[
(68)\qquad \det\sigma N=\operatorname{Tr}\sigma N=0,\quad
\begin{gathered}
\sigma N\text{ non nul en tte fibre.}\\
\text{i.e. }N\text{ non nul en tte fibre}
\end{gathered}
\]
Or la \deleted{conditions} \inserted{troisième de ces} (68) résulte des autres, car si on avait $N=0$ sur une fibre, on pourrait supposer $S$ spectre d'un corps, $N=0$, i.e. $\rho=-\sigma$, donc \deleted{\illegible{1}} on aurait à la fois
\[
\sigma^2+1=0\qquad \sigma^2+\sigma+1=0
\qquad\text{d'où }\sigma=0,\text{ absurde.}
\]
Je dis aussi qu'en aucune fibre on n'a $N=1$, sinon on aurait en cette fibre \deleted{$\rho=1-\sigma$} $\sigma=1-\rho$ et en substituant dans $\sigma^2+1=0$ on trouve \deleted{\illegible{1}}
\[
(1-\rho)^2+1=0\quad\text{i.e.}\quad \rho^2-2\rho+2=0,
\]
\mspage{336}

\scanpage{338}

comparant avec $\rho^2-\rho+1=0$, on trouve \deleted{$\unsure{\rho-1}$} par différence $\rho-1=0$ i.e. $\rho=1$, mais la relation $\rho^2-\rho+1=0$ devient $1=0$, absurde.

Donc l'opérateur idempotent $N$ est \deleted{partout} de rang $1$ \inserted{en chaque fibre}, donc \illegible{1} (sauf erreur) \deleted{\unsure{dit}} que $N$ est \deleted{décomposable} (en termes de $\unsure{\mathcal M}\simeq\check F\otimes F$) \inserted{loc\textsuperscript{t}} décomposable sous la forme $x'\otimes x$, avec $\langle x,x'\rangle=1$, et \deleted{$\operatorname{Tr}$}
\[
\operatorname{Tr}N=1,\quad\det N=0
\]
donc la relation $\det(\sigma N)=0$ dans (68) est également automatique, puisque $\det\sigma N=\det\sigma\det N=0$. Il reste la seule relation
\[
\operatorname{Tr}\sigma N=0
\]
\inserted{$\sigma^2+\sigma\rho=\sigma\rho-1$}

( \deleted{que} qui s'écrit aussi
\[
(69)\qquad \operatorname{Tr}\sigma\rho=2\quad(\Longleftrightarrow\operatorname{Tr}\sigma^{-1}\rho=-2)
\]
\unsure{pour} laquelle il faudrait voir si oui ou non elle résulte des autres relations.

Prenons
\[
(70)\qquad
\left\{\begin{aligned}
\rho&=\rho_0=\begin{pmatrix}0&1\\-1&1\end{pmatrix}\\
\sigma&=\underline u\sigma_0\underline u^{-1}
=\begin{pmatrix}\unsure{a}&b\\c&d\end{pmatrix}
\begin{pmatrix}0&-1\\1&0\end{pmatrix}
\begin{pmatrix}d&-b\\-c&a\end{pmatrix}\\
&\text{où }\underline u=\begin{pmatrix}a&b\\c&d\end{pmatrix}\in\unsure{\Gamma\mathrm{Sl}(2)},\text{ i.e. }ad-bc=1
\end{aligned}\right.
\]
On trouve
\mspage{337}

\scanpage{339}
\mspage{470}

\[
(70\mathrm{bis})\qquad
\sigma=\begin{pmatrix}
bd+ac&-(a^2+b^2)\\
c^2+d^2&-(bd+ac)
\end{pmatrix}
\]
d'où
\[
N=\rho+\sigma=\begin{pmatrix}
(bd+ac)&1-(a^2+b^2)\\
-1+(c^2+d^2)&1-(bd+ac)
\end{pmatrix}
\]
\[
N-1=\begin{pmatrix}
(bd+ac)-1&1-(a^2+b^2)\\
-1+(c^2+d^2)&-(bd+ac)
\end{pmatrix}
\quad\begin{gathered}
=-\det(N)N^{-1}\\
\text{si }\det N\text{ \unsure{inv.}}
\end{gathered}
\]
et on trouve
\[
(71)\qquad N^2-N=N(N-1)=-\det(N).1
\]
donc \deleted{la condition}
\[
(72)\qquad N^2=N\deleted{\unsure{\Longleftrightarrow}}\quad\Longleftrightarrow\quad\det N=0\ ;\quad(\Longrightarrow\operatorname{Tr}N=1)
\]
Or on a
\[
\begin{aligned}
\det N={}&bd+ac-(bd\unsure{+}ac)^2\deleted{\unsure{+1}}\\
&\unsure{-|=}\bigl[-1-(a^2+b^2)(c^2+d^2)+a^2+b^2+c^2+d^2\bigr]\\
={}&bd+ac-\deleted{b^2d^2}-\deleted{a^2c^2}-2a\uline{bc}d+1
+\deleted{a^2c^2}+\uline{a^2d^2}+\uline{b^2c^2}+\deleted{\unsure{b^2d^2}}\\
&\hspace{4em}-a^2-b^2-c^2-d^2\\
={}&(bd+ac)-(a^2+b^2+c^2+d^2)+\underset{1}{\underbrace{(ad-bc)^2}}+1\\
={}&2+(bd+ac)-(a^2+b^2+c^2+d^2)
\end{aligned}
\]
donc la condition (71) équivaut à
\[
(73)\qquad a^2+b^2+c^2+d^2=2+(bd+ac),
\]
par ailleurs on a par 70, (70bis)
\[
\unsure{\rho\sigma}=\begin{pmatrix}
c^2+d^2&-(bd+ac)\\
c^2+d^2-(bd+ac)&a^2+b^2-(bd+ac)
\end{pmatrix}
\]
donc
\[
(74)\qquad \operatorname{Tr}\unsure{\sigma\rho}=\operatorname{Tr}\rho\sigma
=a^2+b^2+c^2+d^2-(bd+ac)=2
\]
\mspage{338}
